\chapter{Adversarial Transfer and Early Behavioral Evidence}
\label{chap:simple-games}
\chapterrunning{adversarial transfer}

The simple-game phases serve two purposes in the dissertation. First, they provide a low-cost
environment in which iterative code evolution can be observed over many epochs without the
computational burden of large benchmark routing suites. Second, they make it possible to inspect
the relationship between code novelty and \emph{behavioral} change before moving to harder
optimization domains. The results from these phases do not establish routing-quality claims, but
they do establish whether the loop can produce transfer-relevant adaptation at all.

\section{Transfer Setting}
\label{sec:games-setting}

The retained simple-game evidence combines two artifacts: the official transfer suite and the
phase-4 novelty adjudication. The transfer suite spans ten replicated runs, three transfer
environments, and 3000 total epochs. The environments are:

\begin{enumerate}[leftmargin=2em]
\item pursuit/evasion,
\item resource collection/denial, and
\item territory control.
\end{enumerate}

The purpose of the suite is not to reopen the earlier curriculum search, but to evaluate whether
the evolved agents remain effective when the environment changes and when performance is measured
on held-out opponent panels. The main quantitative outcomes are hold-out win rate and hold-out
score margin.

The novelty adjudication provides the complementary qualitative control. It reviews the strongest
novelty spikes from three transfer-era curriculum recipes:

\begin{enumerate}[leftmargin=2em]
\item \emph{rotating-opponents holdout endpoint},
\item \emph{rotating + nemesis + novelty + replay}, and
\item \emph{rotating + replay-aware selection}.
\end{enumerate}

Each reviewed spike is conservatively classified as functional adaptation, mixed or local
adjustment, or churn/failed adaptation. This matters because a transfer claim in the early phases
should refer to actual strategy change, not only to different code text.

\section{Manual Novelty Adjudication}
\label{sec:games-novelty}

Table~\ref{tab:games-novelty} summarizes the manual calls for the top novelty spikes in each
recipe.

\begin{table}[t]
\centering
\small
\setlength{\tabcolsep}{6pt}
\begin{tabular}{
  >{\raggedright\arraybackslash}p{2.35in}
  r
  r
  r
}
\toprule
Recipe & Functional & Mixed/local & Churn/failed \\
\midrule
\emph{Rotating-opponents holdout endpoint} & 2/10 & 2/10 & 6/10 \\
\emph{Rotating + nemesis + novelty + replay} & 5/10 & 2/10 & 3/10 \\
\emph{Rotating + replay-aware selection} & 5/10 & 2/10 & 3/10 \\
\bottomrule
\end{tabular}
\caption{Manual adjudication of the strongest novelty spikes in the phase-4 transfer-era recipes.
The heavier replay-aware recipes show more genuine functional adaptations than the simple rotating
baseline, but substantial churn remains even in the stronger recipes.}
\label{tab:games-novelty}
\end{table}

Three points follow from Table~\ref{tab:games-novelty}. First, the baseline rotating-opponent
recipe is dominated by churn or failed adaptation among its largest code changes. Second, the
heavier replay-aware recipes do improve the ratio of useful to useless novelty spikes, which is
evidence that added selection pressure can matter. Third, the stronger recipes still do not justify
a claim of broad strategic invention. Even when they produce functional improvements, those
improvements are uneven across environments and remain mixed with several large but unhelpful code
changes.

\section{Held-Out Transfer Performance}
\label{sec:games-transfer}

The official cross-family synthesis reports the family-level hold-out performance shown in
Figure~\ref{fig:games-transfer}. The numbers are aggregated at the environment level because the
simple-game role in the dissertation is to establish transfer-capable adaptation, not to defend one
specific curriculum recipe as universally optimal.

\begin{figure}[t]
\centering
\begin{tikzpicture}
\begin{groupplot}[
  group style={group size=2 by 1, horizontal sep=1.2cm},
  width=0.38\textwidth,
  height=0.24\textheight,
  ymin=0,
  enlarge x limits=0.18,
  axis line style={black},
  tick style={black},
  xtick=data,
  x tick label style={font=\scriptsize, rotate=22, anchor=east},
  yticklabel style={
    font=\footnotesize,
    /pgf/number format/fixed,
    /pgf/number format/precision=2
  },
  every axis title/.style={font=\small},
  title style={at={(0.5,1.03)}, anchor=south},
  point meta=y,
  nodes near coords={\pgfmathprintnumber[fixed,precision=3]{\pgfplotspointmeta}},
  every node near coord/.append style={font=\tiny, anchor=south, yshift=1pt},
]
\nextgroupplot[
  title={Hold-out win rate},
  ymax=1.05,
  ylabel={Higher is better},
  symbolic x coords={Pursuit,Resource,Territory},
]
\addplot+[ybar, fill=black!20, bar width=10pt] coordinates {
  (Pursuit,0.7000)
  (Resource,0.4840)
  (Territory,0.9467)
};

\nextgroupplot[
  title={Hold-out score margin},
  ymax=36,
  ylabel={},
  symbolic x coords={Pursuit,Resource,Territory},
]
\addplot+[ybar, fill=black!45, bar width=10pt] coordinates {
  (Pursuit,3.702)
  (Resource,1.348)
  (Territory,34.170)
};
\end{groupplot}
\end{tikzpicture}
\caption{Family-level held-out outcomes for the simple-game transfer suite. Territory control
shows the strongest positive transfer, pursuit/evasion remains moderately positive, and
resource-collection/denial is the weakest environment.}
\label{fig:games-transfer}
\end{figure}

The three environments differ materially. Territory control shows the clearest success signal, with
hold-out win rate $0.9467$ and mean score margin $34.17$. Pursuit/evasion is weaker but still
positive, with win rate $0.70$ and margin $3.702$. Resource-collection/denial is the most difficult
transfer environment, with win rate only $0.484$ despite a positive mean margin of $1.348$.

This pattern is consistent with the novelty audit. The most convincing functional adaptations were
concentrated in the territorial environment, while the other environments showed a larger share of
local retuning and churn. The simple-game evidence therefore supports a bounded positive claim:
iterative code evolution can produce real transfer-relevant behavioral change, but that change is
domain-dependent and cannot be inferred from novelty alone.

\section{Interpretation}
\label{sec:games-interpretation}

The simple-game phases support the dissertation in three ways.

First, they provide direct evidence for Hypothesis~H1 from
Section~\ref{sec:intro-hypotheses}. The loop does generate executable diversity, and some of that
diversity corresponds to useful held-out adaptation rather than to random code motion.

Second, they establish an important caution that remains true in later chapters: lexical novelty is
not a reliable proxy for functional gain. Without the manual adjudication pass, the early project
could easily have overstated the meaning of large code changes.

Third, the game phases motivate the move to routing benchmarks. Once the possibility of adaptation
is established in a low-cost environment, the harder question becomes whether the same machinery
survives stronger baselines, clearer notions of feasibility, and more standardized benchmark
evaluation. The next chapter addresses that question directly through replicated TSP, ATSP, and
CVRP studies.
